\documentclass[a4paper,10pt]{article}
\usepackage[utf8]{inputenc}

%Preámbulos de las figuras

%--------------------------------------------
\usepackage{polynom}
%--------------------------------------------
\usepackage{nicematrix,tikz}
\usetikzlibrary{fit}
%--------------------------------------------
\usepackage[table, svgnames]{xcolor}
\usepackage{hhline, booktabs}
%--------------------------------------------
\usepackage{amsmath,gensymb}
\usetikzlibrary{matrix,positioning,calc}
\definecolor{COLOR1}{HTML}{6B99B9}
\definecolor{COLOR2}{HTML}{DD9CC4}
\tikzset{
    M/.style={
        matrix of math nodes,
        nodes={draw=none,anchor=center,minimum size=1.75em},
        nodes in empty cells,
    }
}
%Drawing code nesting in definition with parameters
\def\Det[color #1 size #2 line #3]#4#5{
    \draw[line width=#3,draw=#1, opacity=0.7]
    let \p1 = ($(#5)-(#4)$), % To access cartesian coordinates x, and y.
        \n2 = {atan2(\x1,\y1)}, % to get the angle of the line betwen #5 to #4
    in
    (#5.center)
        ++(90-\n2+90:#2/2)
        arc (90-\n2+90:90-\n2+90-180:#2/2)
        -- ($(#4.center)+(90-\n2-90:#2/2)$)
        arc (90-\n2-90:90-\n2-90-180:#2/2)
        -- cycle;
}
%--------------------------------------------
\usepackage{amssymb}
\newcommand{\ff}{\mathtt{f}}
%--------------------------------------------
\newcolumntype{A}{>{\columncolor{red!20}}c}
\newcolumntype{B}{>{\columncolor{blue!20}}c}

%%------------------------------------------------------------------------


\pagestyle{empty}
\begin{document}
\begin{equation*}
\boldsymbol{\mathcal{F}}^{i} =
\begin{bNiceMatrix}[margin,nullify-dots]
    \ff_{1\,1} & \Cdots & \ff_{1\,\varpi} & \ff_{1\,(\varpi+1)} & \Cdots & \ff_{1\,\eta} \\
    \Vdots & \Ddots & \Vdots & \Vdots & & \Vdots \\
    \ff_{\varpi\,1} & \Cdots & \ff_{\varpi\,\varpi} & \ff_{\varpi\,(\varpi+1)} & \Cdots & \ff_{\varpi\,\eta} \\
    \noalign{\kern2mm}
    \ff_{(\varpi+1)\,1} & \Cdots & \ff_{(\varpi+1)\,\varpi} & \ff_{(\varpi+1)\,(\varpi+1)} & \Cdots & \ff_{(\varpi+1)\,\eta} \\
    \Vdots & & \Vdots & \Vdots & \Ddots & \Vdots \\
    \ff_{\eta\,1} & \Cdots & \ff_{\eta\,\varpi} & \ff_{\eta\,(\varpi+1)} & \Cdots & \ff_{\eta\,\eta}
\CodeAfter
  \SubMatrix[{1-1}{3-3}][slim,extra-height=-1mm,xshift=1mm]
  \OverBrace{1-1}{3-3}{\scriptstyle\mathcal{F}_N^i}[shorten,yshift=1mm]
\end{bNiceMatrix}
\end{equation*}
\end{document}